\documentclass[12pt,twoside]{article}

\usepackage{weiiszablon}

\author{Tomasz Siedlecki, Szymon Jabłoński}

% np. EF-123456, EN-654321, ...
\studentID{173707, 173625}

\title{Metody automatycznej poprawy jakości danych z czujników IoE i ruchu sieciowego}
\titleEN{Methods for automatic enhancement of data quality from IoE sensors and networ traffic}


%%% wybierz rodzaj pracy wpisując jeden z poniższych numerów: ...
% 1 = inżynierska	% BSc
% 2 = magisterska	% MSc
% 3 = doktorska		% PhD
% 4 = praca inżynierska
%%% na miejsce zera w linijce poniżej
\newcommand{\rodzajPracyNo}{0}


%%% promotor
\supervisor{dr inż. Mariusz Bolanowski}

\supervisorEN{Mariusz Bolanowski, PhD}

\abstract{Treść streszczenia po polsku}
\abstractEN{Treść streszczenia po angielsku}

\keywords{max. 5 słów kluczowych w 2 wierszach, oddzielanych przecinkami}
\keywordsEN{max. five keywords in English}


\begin{document}

% strona tytułowa
\maketitle

\blankpage

% spis treści
\tableofcontents

\clearpage
\blankpage

\section{Wstęp}

Dynamiczny rozwój współczesnych środowisk sieciowych, w tym systemów zintegrowanych w ramach koncepcji Internetu Wszechrzeczy (ang. \emph{Internet of Everything} -- IoE), powoduje wytwarzanie ogromych ilości informacji telemetrycznych i o ruchu sieciowym. 
W obszarze cyberbezpieczeństwa kluczowym wyzwaniem staje się skuteczna identyfikacja anomalii i incydentów naruszenia bezpieczeństwa przy użyciu algorytmów uczenia maszynowego.
Możliwość rozpoznawania przez nie wzorców pozwala na wysoką skuteczność w wykrywaniu schematów behawioralnych, w których brakuje jasnej logiki i wyrażnych sygnatur, pozwalających na identyfikację przez tradycyjne systemu.
Efektywność tych modeli jest jednak bezpośrednio uzależniona od zapewnionych danych wejściowych. 
Surowe dane, pochodzące bezpośrednio z sensorów, czy też logów sieciowych, często zawierają w sobie rekordy z anomaliami pomiarowymi, brakującymi wartościami, czy też o niskiej użyteczności bez połączenia z szerszym kontekstem ich powstania.
W związku z tym krytycznego znaczenia nabiera wieloetapowy proces przetwarzania (ang. \emph{data preprocessing pipeline}), którego zadaniem jest transformacja surowych danych do postaci kompletnej i przystosowanej do przetwarzania przez modele.

Niniejsze opracowanie stanowi raport z projektu, mającego za cel zglębienie technik przetwarzania danych w celu poprawnienia rezultatów detekcji anomalii.
Dotyczy on zarówno danych uzyskanych z czujników \texttt{IoE}, jak i zapisów ruchu sieciowego.

% część o czujnikach

% część o ruchu sieciowym
W części projektu dotyczącej danych z ruchu sieciowego użyty został zestaw danych \emph{UNSW\_NB15} \cite{bib:datasetsieciowy}.
Jest to zbiór powstały w Australijskim Centrum dla Cyberbezpieczeństwa (ACCS), stanowi on hybrydowy zbiór zawierający normalny ruch sieciowy, zmieszany z syntetycznymi zachowaniami odwzorowującymi ataki.
Główny punkt ciężkości badań ulokowano w analizie wpływu poszczególnych etapów przygotowania danych na ostateczną architekturę matematyczną modeli uczenia maszynowego. 
Laboratoryjne zestawy danych mają tendencję do monopolizowania uwagi klasyfikatorów przez statyczne zmienne nagłówkowe (takie jak unikalne sygnatury systemów operacyjnych czy sztucznie powtarzalne liczniki pakietów), co w warunkach rzeczywistych prowadzi do całkowitej ślepoty algorytmów. 
W ramach projektu zaimplementowano potok danych, którego kluczowym elementem była transformacja strumienia pakietów do postaci logicznych, dwukierunkowych sesji sieciowych, a następnie poddanie ich rygorystycznemu skalowaniu i kodowaniu kategorycznemu.

W celu weryfikacji wpływu operacji preprocesingu na zdolności detekcyjne, w pracy zaimplementowano i porównano trzy odmienne paradygmaty obliczeniowe: histogramowe wzmacnianie gradientu, las izolacyjny oraz regresję logistyczną z regularyzacją $L_1$. 
Badania przeprowadzono eksperymentalnie na różnych etapach zaawansowania potoku danych. 
Pozwoliło to na matematyczną ocenę, jak operacje takie jak agregacja sesyjna, uzupełnianie braków czy standaryzacja cech wpływają na stabilność wag modeli liniowych oraz na rozkład struktur decyzyjnych w modelach drzewiastych. 
Za pomocą analizy ważności cech (zarówno metodą permutacji, jak i eliminacji cech Lasso) udowodniono, że prawidłowo zaprojektowany preprocesing wymusza na algorytmach porzucenie statycznych skrótów na rzecz analizy rzeczywistych paramertów sesji.

\section{Architektura potoku przetwarzania danych sieciowych}

Wykrywanie anomalii w rzeczywistych warunkach operacyjnych wymaga eliminacji podatności modeli uczenia maszynowego na błędy wynikające ze specyfiki danych testowych. 
W ramach realizowanego projektu zaprojektowano i zaimplementowano modularny potok przetwarzania danych (ang. \emph{data pipeline}), którego nadrzędnym celem była transformacja surowych rekordów sieciowych do postaci niezależnej od konfiguracji środowiska laboratoryjnego. 
Całość algorytmów preprocesingu została zaimplementowana w języku Python z wykorzystaniem bibliotek \texttt{pandas}, \texttt{numpy} oraz \texttt{scikit-learn}. 
Proces przygotowania danych podzielono na sześć sekwencyjnych etapów, sterowanych z poziomu głównego skryptu wykonawczego.

\subsection{Strukturyzacja i przekodowanie danych surowych}

W pierwszym etapie prac (\texttt{step\_1\_add\_column\_names.py}) dokonano strukturyzacji surowych plików tekstowych formatu CSV. 
Ze względu na brak wbudowanych nagłówków w plikach źródłowych, zaimplementowano procedurę automatycznego mapowania etykiet kolumn na podstawie załączonego w zestawie danych pliku z nazwami kolumn.
Ponadto pliki są domyślnie zakodowane w formacie \texttt{Windows 1252}, zostały one przekodowane do standardu \texttt{UTF-8}.
Operacja ta zapewniła integralność schematu danych przed dalszymi przekształceniami i pozwoliła na operacje na nazwach kolumn, ułatwiając prace z nimi.

\subsection{Oczyszczanie i naprawa danych wejściowych}

W drugim etapie (\texttt{step\_2\_fix\_data\_errors.py}) przeprowadzono czyszczenie danych oraz naprawę błędów typowania. 
Implementacja tego modułu była kluczowa z punktu widzenia stabilności modeli matematycznych. 
W ramach walidacji wykonano między innymi następujące operacje:
\begin{itemize}
	\item Sprawdzono poprawność strukturalną adresów IP (źródłowych \emph{srcip} oraz docelowych \emph{dstip}) przy użyciu walidatorów biblioteki \texttt{ipaddress}.
	\item Dokonano konwersji i filtracji portów protokołów warstwy transportowej (\emph{sport}, \emph{dsport}), odrzucając rekordy wykraczające poza zakres od $0$ do $65535$.
	\item Zaimplementowano mechanizm uzupełniania brakujących informacji o usługach (\emph{service}) na podstawie mapowania znanych portów aplikacji (np. porty 80, 443 zmapowano jako usługę HTTP, port 53 jako DNS).
	\item Przeprowadzono filtrację stanów połączeń (\emph{state}), ograniczając je wyłącznie do skończonego zbioru dopuszczalnych flag sieciowych (np. \texttt{FIN}, \texttt{INT}, \texttt{CON}).
\end{itemize}
Po zakończeniu walidacji dokonano jawnej eliminacji wierszy zawierających niekompletne informacje (operacja \texttt{dropna}), co zabezpieczyło algorytmy klasyfikacyjne przed błędami braku danych.

\subsection{Agregacja sesyjna i behawioralna transformacja strumienia}

Kluczowym punktem inżynierii danych, determinującym odporność modeli na warunki rzeczywiste, był krok trzeci (\texttt{step\_3\_group\_by\_sessions.py}). 
Tradycyjne podejście pakietowe jest podatne na zjawisko wycieku danych laboratoryjnych, ponieważ modele uczą się statycznych cech nagłówków pojedynczych ramek. 
W celu wymuszenia analizy parametrów dotyczących całości sesji, zaimplementowano algorytm agregacji sesyjnej.

Pojedyncze rekordy pakietowe zostały zgrupowane w logiczne połączenia przy użyciu klucza grupowania o formacie:
\begin{equation}
	K_{sesji} = \langle \text{srcip}, \text{sport}, \text{dstip}, \text{dsport}, \text{proto} \rangle
\end{equation}

Użyte zostały nazwy cech ze zbioru danych, gdzie:
\begin{itemize}
\item \texttt{srcip} - adres IP źródła,
\item \texttt{sport} - port źródłowy,
\item \texttt{dstip} - adres IP celu,
\item \texttt{dsport} - port docelowy,
\item \texttt{proto} - protokół przesyłanego pakietu,
\end{itemize}

Dla tak zdefiniowanych okien konwersacji wprowadzono zróżnicowaną logikę agregacji parametrów fizycznych i statystycznych:
\begin{itemize}
	\item Cechy wolumetryczne i czasowe, takie jak czas trwania połączenia (\emph{dur}), liczba przesłanych bajtów (\emph{sbytes}, \emph{dbytes}) czy utracone pakiety (\emph{sloss}, \emph{dloss}), zostały zsumowane w obrębie całej sesji.
	\item Parametry dynamiczne, w tym średnie obciążenie sieciowe (\emph{Sload}, \emph{Dload}) oraz fluktuacje opóźnień (tzw. \emph{jitter} -- \emph{Sjit}, \emph{Djit}), zostały przeliczone jako wartości średnie z całego czasu trwania sesji.
	\item Zmienne kategoryczne, takie jak stan połączenia (\emph{state}) czy typ usługi (\emph{service}), przetransformowano za pomocą wyznaczenia wartości najczęstszej w danym oknie czasowym, wyjątkiem były kategorie ataków, które zawsze są zapisywane, nawet jeżeli dominują pakiety normalne.
\end{itemize}
Dzięki tej transformacji zredukowano liczbę rekordów, jednocześnie przekształcając dane z formy izolowanych impulsów sieciowych w ciągłe profile behawioralne.

\subsection{Skalowanie nieliniowe, standaryzacja i kodowanie zmiennych}

W kroku czwartym (\texttt{step\_4\_scale\_standardise.py}) rozwiązano problem skrajnej asymetrii rozkładów cech sieciowych. 
Zmienne takie jak liczba przesłanych bajtów czy czas trwania charakteryzują się skrajnymi wachaniami wartości, co paraliżuje działanie modeli liniowych (np. regresji logistycznej). 
W celu normalizacji rozkładów zastosowano transformację logarytmiczną typu $\log(x + 1)$:
\begin{equation}
	x^* = \ln(x + 1)
\end{equation}
Operację tę wykonano dla cech: \emph{sbytes}, \emph{dbytes}, \emph{Sload}, \emph{Dload}, \emph{dur}, \emph{Sjit}, \emph{Djit} oraz opóźnień TCP (\emph{tcprtt}, \emph{synack}, \emph{ackdat}). 
Pozostałe cechy numeryczne poddano klasycznej standaryzacji przy użyciu algorytmu \texttt{StandardScaler}, sprowadzając ich rozkłady do średniej równej zero oraz odchylenia standardowego równego jeden.

\subsection{Łączenie danych i konwersja cech kategorycznych}

W kroku piątym (\texttt{step\_5\_merge.py}) złączono przetworzone pliki cząstkowe w jeden zbiorczy zbiór danych.
W początkowym zbiorze dane są podzielone na 4 pliki, z maksymalną wielkością ok. 700 tysiecy rekordów.
Ich łączna wielkość to około 2,5 miliona punktów danych.

W ostatnim, szóstym kroku (\texttt{step\_6\_category\_encoding.py}) tekstowe cechy kategoryczne (\emph{proto}, \emph{service}, \emph{state}) zostały zakodowane numerycznie z jednoczesnym wyeksportowaniem słowników mapowania do plików formatu JSON. 
Proces ten został dokonany dopiero po złączeniu plików w celu uniknięcia niespójności w kodowaniu cech w poszczególnych plikach. 

\section{Wyniki eksperymentów i analiza porównawcza dla danych sieciowych}

W celu oceny wpływu zaprojektowanego procesu przetwarzania danych na skuteczność wykrywania anomalii sieciowych, przeprowadzono testy porównawcze na trzech etapach:
\begin{itemize}
\item dla danych surowych z dodanymi etykietami
\item dla danych po oczyszczeniu i naprawie
\item dla danych zgrupowanych
\end{itemize}
We wszystkich przypadkach dane zostały również zgrupowane.

W eksperymentach zestawiono działanie dwóch algorytmów nadzorowanych: klasyfikatora histogramowego wzmacniania gradientu (\emph{HistGradientBoosting}) oraz regresji logistycznej z regularyzacją $L_1$.
Zastosowany został również jeden algorytm nie-nadzorowany - las izolacyjny.

\subsection{Analiza wyników dla danych surowych}

W pierwszej fazie modele zostały uruchomione na danych surowych, gdzie wejściem były pojedyncze pakiety. 
Zbiorcze zestawienie uzyskanych metryk klasyfikacji przedstawia tabela~\ref{Tab:WynikiSurowe}.

\begin{table}[ht]
\caption{Metryki klasyfikacji anomalii dla surowych danych pakietowych}
\centering		
	\begin{tabular}{|l|c|c|c|c|}	
		\hline
		Algorytm & Precision & Recall & F1-Score & ROC-AUC \\
		\hline
		HistGradientBoosting (normalny) & 0,99 & 1,00 & 1,00 & 0,9997 \\
		HistGradientBoosting (anomalie) & 0,98 & 0,96 & 0,97 & 0,9997 \\
		Regresja logistyczna (normalny) $L_1$ & 1,00 & 0,99 & 0,99 & 0,9989 \\
		Regresja logistyczna (anomalie) $L_1$ & 0,93 & 0,97 & 0,95 & 0,9989 \\
		Las izolacyjny (normalny) & 0,91 & 0,91 & 0,91 &  \\
		Las izolacyjny (anomalie) & 0,37 & 0,37 & 0,37 & \\
		\hline
	\end{tabular}	
\label{Tab:WynikiSurowe}
\end{table}

Oba modele nadzorowane osiągnęły na tym etapie bardzo wysokie wskaźniki skuteczności. 
Takie wyniki wydają się być wyśmienitymi, jednak są one wyjaśniane przez rozkład wag cech algorytmu, porównany w późniejszej części raportu.
Model lasu izolacyjnego osiągnął niskie wyniki dla wykrywania anomalii i ogólnie niższe wyniki niż algorytmy nadzorowane.

\subsection{Analiza wyników dla danych po oczyszczeniu i naprawie danych}

W drugiej fazie modele zostały uruchomione na danych oczyszczonych i naprawionych przez drugi krok przetwarzania, gdzie wejściem dalej były pojedyncze pakiety. 
Zbiorcze zestawienie uzyskanych metryk klasyfikacji przedstawia tabela~\ref{Tab:Wyniki2}.

\begin{table}[ht]
\caption{Metryki klasyfikacji anomalii dla oczyszczonych danych pakietowych}
\centering		
	\begin{tabular}{|l|c|c|c|c|}	
		\hline
		Algorytm & Precision & Recall & F1-Score & ROC-AUC \\
		\hline
		HistGradientBoosting (normalny) & 0,99 & 1,00 & 1,00 & 0,9997 \\
		HistGradientBoosting (anomalie) & 0,98 & 0,96 & 0,97 & 0,9997 \\
		Regresja logistyczna (normalny) $L_1$ & 1,00 & 0,99 & 0,99 & 0,9990 \\
		Regresja logistyczna (anomalie) $L_1$ & 0,93 & 0,97 & 0,95 & 0,9990 \\
		Las izolacyjny (normalny) & 0,90 & 0,90 & 0,90 &  \\
		Las izolacyjny (anomalie) & 0,34 & 0,33 & 0,33 & \\
		\hline
	\end{tabular}	
\label{Tab:Wyniki2}
\end{table}

W przypadku algorytmów nadzorowanych zmiany po naprawieniu danych są niewielkie.
Ciekawa zależność ukzała się jednak w wypadku lasu izolacyjnego, którego zdolność do poprawnek klasyfikacji pakietów spadła.

\subsection{Analiza wyników dla danych po agregacji sesyjnej}

W trzeciej fazie badania przeprowadzono na danych, które przeszły agregację, jeszcze bez modyfikacji zakresów i skalowania wartości. 
Metryki dla danych przetworzonych zebrano w tabeli~\ref{Tab:Wyniki3}.

\begin{table}[ht]
\caption{Metryki klasyfikacji anomalii dla danych zagregowanych}
\centering		
	\begin{tabular}{|l|c|c|c|c|}	
		\hline
		Algorytm & Precision & Recall & F1-Score & ROC-AUC \\
		\hline
		HistGradientBoosting (normalny) & 0,99 & 0,99 & 0,99 & 0,9997 \\
		HistGradientBoosting (anomalie) & 0,76 & 0,72 & 0,74 & 0,9997 \\
		Regresja logistyczna (normalny) $L_1$ & 1,00 & 0,98 & 0,99 & 0,9950 \\
		Regresja logistyczna (anomalie) $L_1$ & 0,67 & 1.00 & 0,90 & 0,9950 \\
		Las izolacyjny (normalny) & 0,98 & 0,98 & 0,98 &  \\
		Las izolacyjny (anomalie) & 0,41 & 0,41 & 0,41 & \\
		\hline
	\end{tabular}	
\label{Tab:Wyniki3}
\end{table}

Agregacja sesyjna diametralnie zmienia wyniki otrzymane przez wszystkie modele.
Algorytmy nadzorowane doświadczają tu znacznego spadku dokładności wykrywania podejrzanych zachowań.
Las izolacyjny, odwrotnie, zwiększa swoją efektywność, choć jest ona wciąż poniżej poziomu innych algorytmów.

\subsection{Analiza wyników dla danych po wszystkich krokach}

W czwartej fazie badania przeprowadzono na danych, które przeszły pełną ścieżkę przetwarzania (zamiana pakietów na sesje w kroku 3, logarytmowanie cech wielkościowych w kroku 4 oraz kodowanie zmiennych). 
Metryki dla danych przetworzonych zebrano w tabeli~\ref{Tab:Wyniki4}.

\begin{table}[ht]
\caption{Metryki klasyfikacji anomalii dla danych po wszystkich transformacjach}
\centering		
	\begin{tabular}{|l|c|c|c|c|}	
		\hline
		Algorytm & Precision & Recall & F1-Score & ROC-AUC \\
		\hline
		HistGradientBoosting (normalny) & 0,99 & 0,99 & 0,99 & 0,9955 \\
		HistGradientBoosting (anomalie) & 0,76 & 0,72 & 0,74 & 0,9955 \\
		Regresja logistyczna (normalny) $L_1$ & 1,00 & 0,98 & 0,99 & 0,9938 \\
		Regresja logistyczna (anomalie) $L_1$ & 0,68 & 1.00 & 0,81 & 0,9938 \\
		Las izolacyjny (normalny) & 0,98 & 0,98 & 0,98 & \\
		Las izolacyjny (anomalie) & 0,33 & 0,33 & 0,33 & \\
		\hline
	\end{tabular}	
\label{Tab:Wyniki4}
\end{table}

Po wykonaniu skalowania i standaryzacji wyniki zmieniają się głownie dla lasu izolacyjnego, którego efektywność znacznie spadła.
Algorytmy nadzorowane doświadczają niedużych różnic, co sugeruje niewielki wpływ tego działania.

\section{Wpływ transformacji na dane algorytmy w przypadku danych sieciowych}

\subsection{Histogramowe wzmacnianie gradientu}

\subsubsection{Macierze pomyłek}

W tabeli \ref{tab:cm_histgradient} przedstawione zostały macierze pomyłek dla poszczególnych przeliczeń algorytmu histogramowego wzmacniania gradientu.
Widać na nich niewielki wpływ przejść między 1, a 2 oraz 3, a 4 rodzajami danych.

W pierwszym przypadku spadła w w niewielkim stopniu iość klasyfikacji fałszywie pozytywnych, a bardziej wzrosła ilość fałszywych negatywów.
Trzeba tu zauważyć, że w wyniku usunięcia rekordów błędnych spadła ilość danych, więc tym samym ogólna wydajność klasyfikacji nie wzrosła.

W wypadku standaryzacji i skalowania danych doszło do minimalnej zmiany fałszywych pozytywów, przy spadku fałszywych negatywów.
Tu ilość rekordów pozostała taka sama.

Największa zmiana wydarzyła się przy agregacji danych, w którym to przypadku nastąpiła duża zmiana ich struktury.
Ilość rekordów zmniejszyła się, jednak niepoprawne klasyfikacje wzrosły ilościowo.
Oznacza to że po takich transformacjach model gorzej klasyfikuje dane.

\begin{table}[ht]
\caption{Macierze pomyłek dla klasyfikatora HistGradientBoosting}
\centering
\renewcommand{\arraystretch}{1.3}
\begin{tabular}{|l|c|}
\hline
\textbf{Rodzaj danych} & \textbf{Macierz pomyłek} \\
\hline
1. Dane surowe & $\begin{pmatrix} 663787 & 1843 \\ 3457 & 92928 \end{pmatrix}$ \\
\hline
2. Dane oczyszczone & $\begin{pmatrix} 626487 & 1828 \\ 3566 & 91655 \end{pmatrix}$ \\
\hline
3. Dane zgrupowane& $\begin{pmatrix} 532757 & 4597 \\ 5651 & 14314 \end{pmatrix}$ \\
\hline
4. Dane przetworzone & $\begin{pmatrix} 532751 & 4603 \\ 5526 & 14439 \end{pmatrix}$ \\
\hline
\end{tabular}
\label{tab:cm_histgradient}
\end{table}

\subsubsection{Porównanie wag cech}

Przy obliczeniach algorytmu uzyskana została lista wag przypisanych danym cechom ze zbioru danych.
Jest to bardzo ważna informacja, ponieważ pokazuje ona, które dane są uwzględniane najbardziej przy podejmowaniu decyzji o klasyfikacji.
Ich wartości dla kolejnych wersji modelu zostały przedstawione w tabeli \ref{tab:feature_importances_comparison};
Zawarte w tych listach dane zapewniają ważny dla oceny rzeczywistej oceny klasyfikatora kontekst.
Cecha \texttt{sttl}, oznaczająca wartość TTL pakietu od źródła do celu, stanowi bardzo silną sygnaturę syntetyczności danych o atakach.
W bardzo dużej części rekordów reprezentujących ataki wartość ta jest ustawiona na 255.
Model przypisał jej więc bardzo wysoką wagę, ograniczając wpływ innych cech, powodując tzw. \texttt{overfitting}.
Po zgrupowaniu, mimo pozostawieniu maksymalnej wartości tej cechy w ramach grupy, jej waga spadła znacząco.
Większą część ogółu branych pod uwagę cech zaczęły stanowić cechy niepowiązane.
Oznacza to, że mimo spadku poprawności klasyfikacji, modele oparte a dane zagregowane klasyfikują w o wiele większym stopniu w oparciu o parametry behawieralne sesji, zamiast łatwych do podmiany sygnatur. 

\begin{table}[htbp]
\caption{Porównanie wartości ważności cech (Feature Importances) dla klasyfikatora HistGradientBoosting dla czterech zestawów danych}
\centering
\begin{tabular}{|l|c|c|c|c|}
\hline
\textbf{Nazwa cechy} & \textbf{1. Surowe} & \textbf{2. Oczyszczone} & \textbf{3. Zgrupowane} & \textbf{4. Pełne} \\
\hline
\texttt{sttl} & 0,1362 & 0,1549 & 0,0546 & 0,0531 \\
\texttt{ct\_state\_ttl} & 0,0092 & 0,0107 & 0,0010 & 0,0016 \\
\texttt{dsport} & 0,0043 & 0,0040 & 0,0020 & 0,0021 \\
\texttt{service} & 0,0020 & 0,0017 & 0,0004 & 0,0004 \\
\texttt{dbytes} & 0,0017 & 0,0018 & 0,0000 & 0,0001 \\
\texttt{sport} & 0,0010 & 0,0012 & 0,0003 & 0,0003 \\
\texttt{ct\_srv\_dst} & 0,0009 & 0,0008 & 0,0014 & 0,0016 \\
\texttt{smeansz} & 0,0007 & 0,0008 & 0,0006 & 0,0005 \\
\texttt{sbytes} & 0,0005 & 0,0007 & 0,0001 & 0,0002 \\
\texttt{synack} & 0,0005 & 0,0005 & 0,0000 & 0,0000 \\
\texttt{proto} & 0,0005 & 0,0007 & 0,0000 & 0,0000 \\
\texttt{ct\_srv\_src} & 0,0005 & 0,0006 & 0,0005 & 0,0005 \\
\texttt{state} & 0,0004 & 0,0002 & 0,0002 & 0,0006 \\
\texttt{Dload} & 0,0004 & 0,0003 & 0,0000 & 0,0000 \\
\texttt{dmeansz} & 0,0004 & 0,0006 & 0,0003 & 0,0002 \\
\texttt{dloss} & 0,0003 & 0,0002 & 0,0000 & 0,0000 \\
\texttt{ct\_flw\_http\_mthd} & 0,0003 & 0,0000 & 0,0000 & 0,0000 \\
\texttt{ct\_dst\_src\_ltm} & 0,0002 & 0,0004 & 0,0001 & 0,0001 \\
\texttt{dur} & 0,0002 & 0,0006 & 0,0000 & 0,0001 \\
\texttt{Sload} & 0,0002 & 0,0000 & 0,0002 & 0,0000 \\
\texttt{ct\_dst\_sport\_ltm} & 0,0002 & 0,0002 & 0,0000 & 0,0000 \\
\texttt{Dpkts} & 0,0002 & 0,0001 & 0,0000 & 0,0000 \\
\texttt{ct\_src\_ltm} & 0,0001 & 0,0001 & 0,0000 & 0,0000 \\
\texttt{Dintpkt} & 0,0001 & 0,0001 & 0,0001 & 0,0001 \\
\texttt{ct\_src\_dport\_ltm} & 0,0001 & 0,0001 & 0,0000 & 0,0000 \\
\texttt{Djit} & 0,0001 & 0,0001 & 0,0000 & 0,0001 \\
\texttt{ct\_dst\_ltm} & 0,0001 & 0,0001 & 0,0001 & 0,0000 \\
\texttt{dtcpb} & 0,0001 & 0,0001 & 0,0000 & 0,0000 \\
\texttt{Spkts} & 0,0001 & 0,0000 & 0,0000 & 0,0000 \\
\texttt{Sintpkt} & 0,0001 & 0,0000 & 0,0000 & 0,0000 \\
\texttt{sloss} & 0,0001 & 0,0000 & 0,0000 & 0,0000 \\
\texttt{res\_bdy\_len} & 0,0000 & 0,0000 & 0,0000 & 0,0001 \\
\texttt{Sjit} & 0,0000 & 0,0001 & 0,0000 & 0,0000 \\
\texttt{stcpb} & 0,0000 & 0,0000 & 0,0000 & 0,0000 \\
\texttt{is\_ftp\_login} & 0,0000 & 0,0000 & 0,0000 & 0,0000 \\
\texttt{ct\_ftp\_cmd} & 0,0000 & 0,0000 & 0,0000 & 0,0000 \\
\hline
\multicolumn{5}{|c|}{\textit{Cechy wprowadzone / ujawnione wyłącznie w późniejszych krokach}} \\
\hline
\texttt{tcprtt} & - & 0,0001 & 0,0000 & 0,0000 \\
\texttt{trans\_depth} & - & 0,0000 & 0,0000 & 0,0000 \\
\texttt{dttl} & - & 0,0000 & 0,0000 & 0,0000 \\
\texttt{swin} & - & - & - & 0,0000 \\
\texttt{dwin} & - & - & - & 0,0000 \\
\hline
\end{tabular}
\label{tab:feature_importances_comparison}
\end{table}

\subsection{Regresja Logistyczna $L_1$}

\subsubsection{Macierze pomyłek}

W tabeli \ref{tab:cm_logistic_regression} przedstawione zostały macierze pomyłek dla poszczególnych przeliczeń algorytmu regresji logistycznej z regularyzacją $L_1$. 
Podobnie jak w przypadku klasyfikatora histogramowego, widoczny jest minimalny wpływ czyszczenia oraz standaryzacji danych na ostateczną strukturę błędów, przy jednoczesnym drastycznym przemodelowaniu wyników po zastosowaniu agregacji potoku sesyjnego.

Wdrażając grupowanie danych, zaobserwować można skokowy wzrost liczby poprawnie wykrytych anomalii przy jednoczesnym wzroście fałszywych pozytywów. 
Świadczy to o silnym wpływie parametru zbalansowania wag klas (\texttt{class\_weight='balanced'}), który zmusił model do kategoryzacji większości niejednoznacznych sesji jako potencjalnych ataków, minimalizując krytyczny z punktu widzenia bezpieczeństwa współczynnik fałszywych negatywów (FN spadające do poziomu kilkudziesięciu rekordów).

\begin{table}[ht]
\caption{Macierze pomyłek dla klasyfikatora Regresji Logistycznej $L_1$}
\centering
\renewcommand{\arraystretch}{1.3}
\begin{tabular}{|l|c|}
\hline
\textbf{Rodzaj danych} & \textbf{Macierz pomyłek} \\
\hline
1. Dane surowe & $\begin{pmatrix} 658435 & 7195 \\ 2827 & 93558 \end{pmatrix}$ \\
\hline
2. Dane oczyszczone & $\begin{pmatrix} 621368 & 6947 \\ 2869 & 92352 \end{pmatrix}$ \\
\hline
3. Dane zgrupowane & $\begin{pmatrix} 527677 & 9677 \\ 64 & 19901 \end{pmatrix}$ \\
\hline
4. Dane przetworzone & $\begin{pmatrix} 528008 & 9346 \\ 80 & 19885 \end{pmatrix}$ \\
\hline
\end{tabular}
\label{tab:cm_logistic_regression}
\end{table}

\subsubsection{Porównanie współczynników wagowych cech}

W przypadku regresji logistycznej znaczenie danych wejściowych interpretowane jest bezpośrednio poprzez bezwzględne wartości współczynników wagowych (kontrybucji w zmianę logarytmu szans) przypisanych poszczególnym zmiennym w procesie optymalizacji. 
Dla cech kategorycznych poddanych kodowaniu zero-jedynkowemu, takich jak \texttt{proto}, \texttt{state} czy \texttt{service}, w tabeli \ref{tab:feature_weights_regression} zaprezentowano maksymalną bezwzględną wartość współczynnika spośród wszystkich wygenerowanych podklas.

Tak jak w przypadku histogramowego wzmacniania gradientu, zgrupowanie zwiększa udział innych cech sesji, realnie poprawniając efektywność modelu.
Świadczy o tym również bardzo mała ilość fałszywych negatywów.

\begin{table}[htbp]
\caption{Porównanie bezwzględnych wartości wag cech dla klasyfikatora Regresji Logistycznej $L_1$ (Lasso)}
\centering
\begin{tabular}{|l|c|c|c|c|}
\hline
\textbf{Nazwa cechy} & \textbf{1. Surowe} & \textbf{2. Oczyszczone} & \textbf{3. Zgrupowane} & \textbf{4. Pełne} \\
\hline
\texttt{sttl} & 1,0423 & 1,1382 & 0,4032 & 0,3486 \\
\texttt{ct\_state\_ttl} & 1,8838 & 1,8810 & 0,2531 & 0,1614 \\
\texttt{dsport} & 0,0086 & 0,0015 & 0,0403 & 0,1170 \\
\texttt{service} & 0,5779 & 0,6176 & 0,3492 & 0,4043 \\
\texttt{dbytes} & 0,0719 & 0,0810 & 0,0295 & 0,1332 \\
\texttt{sport} & 1,1301 & 1,1493 & 0,6691 & 0,7264 \\
\texttt{ct\_srv\_dst} & 0,1766 & 0,1696 & 0,2112 & 0,2047 \\
\texttt{smeansz} & 0,1086 & 0,0951 & 0,0217 & 0,1141 \\
\texttt{sbytes} & 0,1350 & 0,1411 & 0,0770 & 0,2565 \\
\texttt{synack} & 0,3105 & 0,3174 & 0,2314 & 0,3736 \\
\texttt{proto} & 0,4467 & 0,4478 & 0,3894 & 0,3406 \\
\texttt{ct\_srv\_src} & 0,1805 & 0,1654 & 0,1178 & 0,1107 \\
\texttt{state} & 0,7932 & 0,7955 & 0,5322 & 0,7556 \\
\texttt{Dload} & 0,1677 & 0,2110 & 0,0535 & 0,1801 \\
\texttt{dmeansz} & 0,6613 & 0,7076 & 0,4771 & 0,2562 \\
\texttt{dloss} & 0,0096 & 0,0090 & 0,0058 & 0,0187 \\
\texttt{ct\_flw\_http\_mthd} & 0,0247 & 0,0254 & 0,0230 & 0,0351 \\
\texttt{ct\_dst\_src\_ltm} & 0,1406 & 0,1573 & 0,0652 & 0,0457 \\
\texttt{dur} & 0,0099 & 0,0137 & 0,0446 & 0,2642 \\
\texttt{Sload} & 0,0525 & 0,0520 & 0,0194 & 0,0210 \\
\texttt{ct\_dst\_sport\_ltm} & 0,1508 & 0,1615 & 0,2224 & 0,1736 \\
\texttt{Dpkts} & 0,0208 & 0,0214 & 0,0179 & 0,0270 \\
\texttt{ct\_src\_ltm} & 0,0933 & 0,0903 & 0,0416 & 0,0459 \\
\texttt{Dintpkt} & 0,0050 & 0,0332 & 0,0648 & 0,3251 \\
\texttt{ct\_src\_dport\_ltm} & 0,0195 & 0,0176 & 0,0993 & 0,0789 \\
\texttt{Djit} & 0,0388 & 0,0369 & 0,0360 & 0,0314 \\
\texttt{ct\_dst\_ltm} & 0,0171 & 0,0160 & 0,0053 & 0,0269 \\
\texttt{dtcpb} & 0,0439 & 0,0386 & 0,0550 & 0,0191 \\
\texttt{Spkts} & 0,0176 & 0,0340 & 0,0705 & 0,0917 \\
\texttt{Sintpkt} & 0,0112 & 0,0707 & 0,0369 & 0,0199 \\
\texttt{sloss} & 0,2479 & 0,2673 & 0,1513 & 0,1391 \\
\texttt{res\_bdy\_len} & 0,0763 & 0,0698 & 0,0144 & 0,2602 \\
\texttt{Sjit} & 0,0237 & 0,0234 & 0,0234 & 0,1629 \\
\texttt{stcpb} & 0,0487 & 0,0520 & 0,0578 & 0,0252 \\
\texttt{is\_ftp\_login} & 0,1859 & 0,2503 & 0,1304 & 0,0911 \\
\texttt{ct\_ftp\_cmd} & 0,0441 & 0,0542 & 0,0491 & 0,1003 \\
\hline
\multicolumn{5}{|c|}{\textit{Cechy wprowadzone / ujawnione wyłącznie w późniejszych krokach}} \\
\hline
\texttt{tcprtt} & 0,1676 & 0,1683 & 0,1254 & 0,1107 \\
\texttt{trans\_depth} & 0,1669 & 0,1826 & 0,0790 & 0,1366 \\
\texttt{dttl} & 0,8127 & 0,7909 & 1,5971 & 1,9488 \\
\texttt{swin} & 0,3371 & 0,3363 & 0,3607 & 0,3131 \\
\texttt{dwin} & 0,1997 & 0,1927 & 0,2907 & 0,2212 \\
\hline
\end{tabular}
\label{tab:feature_weights_regression}
\end{table}

\subsection{Las Izolacyjny}

\subsubsection{Macierze pomyłek}

W tabeli \ref{tab:cm_isolation_forest} zestawione zostały macierze pomyłek dla algorytmu Lasu Izolacyjnego na czterech etapach przetwarzania danych. 
Ze względu na nienadzorowany charakter tego klasyfikatora, kluczowym czynnikiem konfiguracji było jawne zdefiniowanie parametru zanieczyszczenia (\emph{contamination}), określającego spodziewany odsetek anomalii w analizowanym zbiorze danych.

\begin{table}[ht]
\caption{Macierze pomyłek dla klasyfikatora Isolation Forest}
\centering
\renewcommand{\arraystretch}{1.3}
\begin{tabular}{|l|c|}
\hline
\textbf{Rodzaj danych} & \textbf{Macierz pomyłek} \\
\hline
1. Dane surowe & $\begin{pmatrix} 2015438 & 203326 \\ 203351 & 117932 \end{pmatrix}$ \\
\hline
2. Dane oczyszczone & $\begin{pmatrix} 1893107 & 201277 \\ 213636 & 103766 \end{pmatrix}$ \\
\hline
3. Dane zgrupowane & $\begin{pmatrix} 1751802 & 39377 \\ 39383 & 27167 \end{pmatrix}$ \\
\hline
4. Dane przetworzone & $\begin{pmatrix} 1746621 & 44558 \\ 44564 & 21986 \end{pmatrix}$ \\
\hline
\end{tabular}
\label{tab:cm_isolation_forest}
\end{table}

\subsubsection{Analiza zachowania modelu nienadzorowanego}

W odróżnieniu od nadzorowanych modeli drzewiastych i liniowych, Isolation Forest izoluje anomalie w oparciu o ich odrębność topologiczną w przestrzeni cech, a nie na podstawie dostarczonych etykiet klas. 
Analiza macierzy pomyłek ujawnia specyficzną symetrię błędów – liczba fałszywych alarmów na każdym z etapów jest niemal identyczna jak liczba niezidentyfikowanych ataków. 
Wynika to bezpośrednio z matematycznej konstrukcji algorytmu: model odcina dokładnie taki ułamek próbki jako anomalie, jaki zdefiniowano w parametrze \texttt{contamination}, co przy niedoskonałej separacji klas skutkuje równomiernym rozkładem pomyłek w obie strony.

Wprowadzenie agregacji sesyjnej zaowocowało drastycznym spadkiem bezwzględnej liczby błędów obu typów oraz wzrostem ogólnej dokładności (\emph{accuracy}) do poziomu 0,96. 
Wynika to ze zmniejszenia gęstości szumu informacyjnego – pojedyncze, surowe pakiety w sieci wykazują ogromną zmienność losową, przez co algorytm nienadzorowany błędnie izolował poprawne pakiety o nietypowym rozmiarze. 
Złączenie ruchu do postaci sesji nadało danym bardziej wyrazistą i zwartą strukturę geometryczną.

Niezwykle interesujące jest pogorszenie wyników klasyfikacji po zastosowaniu pre-processingu, standaryzacji i kodowania \emph{one-hot} (etap 4), gdzie liczba poprawnie wykrytych ataków spadła z 27167 do 21986. 
Zjawisko to wynika ze specyfiki działania Lasu Izolacyjnego. 
Przeprowadzone skalowanie (\emph{StandardScaler} lub \emph{MinMaxScaler}) wyrównało wariancję poszczególnych cech i sprowadziło je do wspólnej skali geometrycznej. 
Spowodowało to, że cechy silnie skorelowane z atakami przestały dominować swoją pierwotną, dużą wartością, a zakodowane zmienne kategorialne sztucznie zagęściły przestrzeń decyzyjną. 
Dla algorytmu szukającego "odizolowanych punktów" transformacja ta realnie utrudniła budowanie losowych podziałów i doprowadziła do częściowego "ukrycia" anomalii pośród ruchu prawidłowego. 
Dowodzi to, że nienadzorowane metody metryczne i podziału przestrzeni często funkcjonują efektywniej na danych nieskalowanych, zachowujących naturalne dla danej domeny proporcje odległości.

\subsection{Porównanie modeli}

Zestawienie wyników klasyfikacji oraz ewolucji struktur decyzyjnych dla modeli Histogramowego Wzmacniania Gradientu, Regresji Logistycznej $L_1$ oraz Lasu Izolacyjnego ujawnia fundamentalne różnice w ich matematycznej wrażliwości na transformacje potoku danych.
Każdy z analizowanych algorytmów reprezentuje odmienne podejście, co przełożyło się na skrajnie różne reakcje na czyszczenie, agregację oraz standaryzację danych.

Na etapach analizy pakietowej (etap 1 i 2), modele nadzorowane (Histogramowe Wzmacnianie Gradientu i Regresja Logistyczna) wykazały silną tendencję do tzw. \emph{overfittingu} (nadmiernego dopasowania) do laboratoryjnych artefaktów środowiska testowego. 
Zjawisko to objawiło się poprzez przypisanie dominujących wag i współczynników zmiennym statycznym, takim jak port źródłowy \texttt{sport} (waga 1,14 w regresji) oraz nagłówkom \texttt{sttl} i \texttt{ct\_state\_ttl}. 
Z punktu widzenia inżynierii bezpieczeństwa sieciowego oznacza to, że klasyfikatory zamiast identyfikować behawioralne wzorce anomalii, nauczyły się rozpoznawać sygnatury konkretnych systemów operacyjnych maszyn generujących ruch w środowisku testowym. 
W tym samym czasie nienadzorowany Las Izolacyjny na danych surowych zmagał się z ogromnym szumem informacyjnym. Wysoka różnorodność pojedynczych pakietów uniemożliwiła precyzyjne odizolowanie punktów, generując duży współczynnik fałszywych alarmów przy niskiej czułości wykrywania rzeczywistych ataków.

Przełomowym krokiem metodologicznym okazała się agregacja danych do postaci dwukierunkowych sesji sieciowych (etap 3). Przekształcenie to wymusiło geometryczną reorganizację przestrzeni cech:
\begin{itemize}
    \item W modelu \textbf{Histogramowego Wzmacniania Gradientu} nastąpiła decentralizacja procesu decyzyjnego – waga zmiennej \texttt{sttl} spadła trzykrotnie (z 0,15 na 0,05), uwalniając masę decyzyjną na rzecz innych cech strukturalnych. 
    Mimo pewnego wzrostu błędów, model zyskał odporność na potencjalne zmiany w syganturach nagłówków.
    \item W \textbf{Regresji Logistycznej} mechanizm regularyzacji $L_1$ (Lasso) zareagował drastycznym zmniejszeniem znaczenia \texttt{ct\_state\_ttl} na rzecz cechy \texttt{dttl} (skok współczynnika do 1,59). 
    Zastosowanie zbalansowanych wag klas (\texttt{class\_weight='balanced'}) zminimalizowało liczbę fałszywych negatywów do krytycznie niskiego poziomu (jedynie 64 przeoczone ataki), kosztem akceptowalnego wzrostu fałszywych alarmów.
    \item Dla \textbf{Lasu Izolacyjnego} konsolidacja ruchu do formy sesji ujednoliciła strukturę geometryczną danych poprawnych. 
    Spowodowało to skokowy wzrost ogólnej dokładności (\emph{accuracy} do 0,96) oraz redukcję błędów obu typów o ponad 80\%, co potwierdza, że nienadzorowane metody podziału przestrzeni wymagają danych opartych o większe zgrupowane zestawy, zamiast pojedynczych punktów.
\end{itemize}

Ostatni etap przetwarzania – pre-processing obejmujący standaryzację, skalowanie oraz pełne kodowanie kategorialne (etap 4) zróżnicował ponowie analizowane algorytmy. 
Dla nieliniowego modelu Histogramowego Wzmacniania Gradientu krok ten okazał się neutralny, drzewa decyzyjne z natury posiadają znaczną odporność na zmiany skali wartości.
Silny efekt odnotowano natomiast w regresji logistycznej: ujednolicenie wariancji znacznie zwiększyły udział cech reprezentujących dynamikę zachowania w sesji.
Cechy określające dynamikę czasową i wolumenową sesji, takie jak \texttt{Dintpkt} (wzrost z 0,06 na 0,32) czy \texttt{dur} (wzrost z 0,04 na 0,26), zyskały kluczowy wpływ na ostateczny klasyfikator. 

Zupełnie odwrotny, negatywny skutek standaryzacji zaobserwowano w przypadku Lasu Izolacyjnego, gdzie liczba poprawnie wykrytych anomalii spadła z 27167 do 21986 dla analizowanego przypadku. 
Przeskalowanie danych zniwelowało naturalne dysproporcje odległości i sztucznie zagęściło wielowymiarową przestrzeń, utrudniając algorytmowi losową izolację unikalnych ścieżek drzew.

Podsumowując, optymalnym wyborem dla systemów detekcji intruzów (IDS) pracujących na danych zagregowanych pozostaje nieliniowy model \textbf{Histogramowego Wzmocnienia Gradientu}, ze względu na stabilność, brak podatności na brak skalowania oraz zrównoważony profil cech. 
\textbf{Regresja Logistyczna} po pre-processingu stanowi dla niego wysoce wydajną, przewidywalną alternatywę skupioną na wykrywaniu ataków, kosztem większej ilości fałszywych pozytywów. 

Z kolei \textbf{Las Izolacyjny}, jako jedyny model nienadzorowany, posiada pewien potencjał dla danych zogranizowanych w sesje.
Mógłby być to dobry wybór dla wykrywania anomalii w dłuższych sesji, zamiast w czasie rzeczywistym.

\subsubsection{Podsumowanie i wnioski końcowe}

% czujniki

% ruch sieciowy

W przypadku danych o ruchu sieciowym rzeprowadzone eksperymenty z użyciem trzech różnorodnych algorytmów uczenia maszynowego pozwoliły wyciągnięcie pewnych obserwacji.
Zastosowanie wielu prób w różnych etapach transformacji pozwoliło ustalić sytuacje w których najlepiej funkcjonują dane podejścia:

\begin{enumerate}
    \item \textbf{Rola grupowania danych:} Przejście z analizy pojedynczych pakietów do analizy dwukierunkowych sesji sieciowych okazało się najważniejszym czynnikiem eliminującym zjawisko przeuczenia (\emph{overfitting}) do sztucznego środowiska laboratoryjnego.
    Konsolidacja ruchu pozwoliła wszystkim modelom na rezygnację z użycia, niestabilnych w rzeczywistym świecie, sygnatur sprzętowych (np. \texttt{sport}, \texttt{ct\_state\_ttl}) na rzecz bardziej użytecznych cech strukturalnych.
    \item \textbf{Skalowanie i standaryzacja:} Badania dowiodły, że transformacje wartości danych (skalowanie, standaryzacja wariancji) nie są zawsze dobrym wyborem i ich stosowanie musi być ściśle dopasowane do natury matematycznej algorytmu. 
    O ile dla modelu \textbf{Histogramowego Wzmacniania Gradientowego} zmiany te nie zmieniały w dużym stopniu wyniku, o tyle dla \textbf{Regresji Logistycznej} stanowił warunek konieczny do prawidłowego ukształtowania liniowej hiperpłaszczyzny rozdzielającej. 
    Z kolei dla \textbf{Lasu Izolacyjnego} zabieg ten okazał się destrukcyjny, sztucznie zagęszczając przestrzeń i utrudniając izolację anomalii.
    \item \textbf{Przewaga modeli nadzorowanych:} Choć autonomiczne podejście Lasu Izolacyjnego wykazało wysoką skuteczność na danych sesyjnych bez dostępu do etykiet klas, to modele nadzorowane zapewniają znacznie większą stabilność i precyzję. 
    W szczególności model \textbf{HistGradientBoosting} charakteryzował się najwyższą odpornością na szum informacyjny, natomiast \textbf{Regresja Logistyczna $L_1$} z mechanizmem bilansowania klas pozwoliła na niemal całkowitą eliminację fałszywych negatywów.
\end{enumerate}

Ostatecznie, uzyskane wyniki potwierdzają sugerują, że kluczem do skutecznej detekcji intruzów w systemach sieciowych jest nie tylko sam wybór wysoce nieliniowego klasyfikatora, ale przede wszystkim odpowiednia architektura potoku przetwarzania danych (\emph{data pipeline}), która przekształca surowy, chaotyczny strumień pakietów w uporządkowaną strukturę sesyjno-behawioralną.





\clearpage

\addcontentsline{toc}{section}{Literatura}

\begin{thebibliography}{4}
\bibitem{bib:datasetsieciowy} https://www.kaggle.com/datasets/mrwellsdavid/unsw-nb15/data?select=UNSW\_NB15\_testing-set.csv. Dostęp 31.05.2026.
\end{thebibliography}

\clearpage

\makesummary

\end{document} 
